\documentclass[11pt,a4paper]{article}

% --------------------------------------------------
% Packages and formatting
% --------------------------------------------------
\usepackage[T1]{fontenc}
\usepackage[utf8]{inputenc}
\usepackage{lmodern}
\usepackage[margin=1in]{geometry}
\usepackage{amsmath,amssymb}
\usepackage{microtype}
\usepackage[hidelinks]{hyperref}
\usepackage{url}

\setlength{\parindent}{1.5em}
\setlength{\parskip}{0.25em}
\emergencystretch=2em

% --------------------------------------------------
% Title information
% --------------------------------------------------
\title{
    A Comparative Study of Classical Numerical Methods,\\
    Physics-Informed Neural Networks, and DeepONets\\
    for Option Pricing and Sensitivity Analysis
}

\author{
    Shivansh Jain\\
    Computer Science and Engineering\\
    LNMIIT
}

\date{September 2026}

\begin{document}

\maketitle

% --------------------------------------------------
% Abstract
% --------------------------------------------------
\begin{abstract}
This project proposes a comparative study of classical numerical
methods, physics-informed neural networks (PINNs), and deep operator
networks (DeepONets) for option valuation and sensitivity analysis.
The study combines controlled Black--Scholes benchmarks with an
empirical case study using Indian equity and derivatives data.
The first stage considers single-asset European options, comparing
Monte Carlo simulation, the Crank--Nicolson finite-difference method,
PINNs, and DeepONets against the analytical solution. Evaluation
will cover pricing and Greek errors, computational cost, training
reliability, repeated valuation, and generalisation across
moneyness, maturity, volatility, and interest rates. Calibration
experiments will examine parameter recovery from synthetic
observations and pricing on held-out market observations.
Model-generated solutions will provide labelled training data
where required, while physics-informed models will use the
governing equation, terminal payoff, and boundary conditions as
training constraints. The second stage considers basket options
on up to five correlated assets, with historical stock returns
informing volatility and correlation inputs. Asset dimension,
correlation structure, and basket weights will also vary.
Monte Carlo with variance reduction and randomised quasi-Monte
Carlo will provide numerical baselines. Analytical geometric
basket values and independent simulation estimates will support
validation. Training, data-generation, and evaluation costs will
be reported separately, and numerical approximation error will
be distinguished from disagreement with observed market prices.
The objective is to identify when neural methods offer useful
computational benefits and when established methods remain
more effective.
\end{abstract}

\noindent
\textbf{Keywords:}
Black--Scholes equation; option pricing; physics-informed neural
networks; DeepONet; Monte Carlo; Greeks; basket options;
parameter calibration.

% --------------------------------------------------
% Introduction
% --------------------------------------------------
\section{Introduction}

Option pricing is an important application of mathematical
modelling in finance. Besides estimating the value of a
contract, pricing methods support risk measurement, hedging,
and calibration to market observations. The Black--Scholes
framework provides a well-established starting point for
these tasks and a controlled setting in which different
computational methods can be evaluated
\cite{black1973,merton1973}.

Classical approaches, including finite differences and
Monte Carlo simulation, remain effective tools for option
valuation. Neural approaches offer alternative ways to
approximate pricing functions. Physics-informed neural
networks (PINNs) incorporate the governing equation into
training, while deep operator networks (DeepONets) are
designed to learn mappings between families of input and
output functions \cite{raissi2019,lu2021}. Their potential
for repeated evaluation and sensitivity analysis motivates
further investigation. However, training cost and reliability
must be considered alongside accuracy, and a neural method
cannot be assumed to outperform an established solver.

This project compares classical methods, PINNs, and
DeepONets in two stages. The first considers single-asset
European options, where the analytical Black--Scholes
solution provides a reference. The second extends the
comparison to basket options on up to five correlated
assets. The assessment goes beyond pricing to include
Greeks, computational efficiency, generalisation, and
parameter calibration. Particular interest lies in whether
a trained model offers useful savings when many related
valuations are required.

Controlled model-generated benchmarks will be complemented
by a case study using Indian equity and derivatives data
from the NSE~\cite{nseData}. This combination allows numerical accuracy
to be examined separately from agreement with market
observations. Historical stock data will also inform the
inputs of the basket experiments.

The objective is to establish the strengths and limitations
of each approach through a reproducible comparison.
Rather than assuming a general advantage for neural
methods, the study will identify the circumstances in
which their accuracy, reliability, and computational
requirements make them a useful alternative.

\section{Classical Pricing Methods: Theory and Implementation}
\label{sec:classical}

This section develops the classical baselines for the comparison with PINNs and DeepONets. The analytical Black--Scholes formula provides reference prices and sensitivities in one dimension. Monte Carlo and randomised quasi-Monte Carlo approximate the risk-neutral expectation, while Crank--Nicolson solves a discretised pricing PDE. The implementation procedures below specify planned experiments; they do not imply that the full comparison has already been completed.

\subsection{Pricing assumptions and terminal simulation}

Consider a European option with strike price \(K\), the contractual exercise price, and current underlying asset price \(S\). Let \(t\) be the current time, \(T\) the expiry time, and \(\tau=T-t\) the time remaining, measured in years. The process \(S_t\) denotes the underlying asset price at time \(t\), with \(S_t=S\) at valuation. The continuously compounded risk-free interest rate \(r\), annualised volatility \(\sigma\), and continuous dividend yield \(q\) are constant; rates are expressed as decimals. Under the risk-neutral measure,
\begin{equation}
 dS_t=(r-q)S_t\,dt+\sigma S_t\,dW_t^Q.
 \label{eq:s3_sde}
\end{equation}
Here \(W_t^Q\) is standard Brownian motion under the risk-neutral probability measure \(Q\), and \(\mathbb E^Q\) denotes expectation under that measure. The initial controlled experiments set \(q=0\). Dividend adjustments, when needed for market observations, will be applied consistently across methods. This treatment is an approximation to dividends and is not an exact model of discrete dividend payments.

Applying It\^o's lemma to \(\log S_t\) and integrating gives
\begin{equation}
 S_T=S\exp\left[\left(r-q-\frac{\sigma^2}{2}\right)\tau
       +\sigma\sqrt{\tau}\,Z\right],\qquad Z\sim\mathcal N(0,1).
 \label{eq:s3_terminal}
\end{equation}
For a terminal payoff function \(g\), the option value \(V(S,t)\) is
\begin{equation}
 V(S,t)=e^{-r\tau}\mathbb E^Q[g(S_T)\mid S_t=S].
 \label{eq:s3_expectation}
\end{equation}
Calls use \(g(S_T)=(S_T-K)^+\), and puts use \(g(S_T)=(K-S_T)^+\), where \(x^+=\max(x,0)\). The analytical formula, simulation methods, and PDE solver target this same model value.

Equation~\eqref{eq:s3_terminal} samples the terminal distribution exactly under the assumptions. European vanilla and basket payoffs depend only on terminal prices, so full time paths are unnecessary for their valuation. Direct terminal sampling avoids time-discretisation error and reduces storage. Full-path simulations can still be used to validate the stochastic model, but are unnecessary for these terminal-payoff pricing experiments.

\subsection{Analytical Black--Scholes benchmark}

\subsubsection{Prices and limiting cases}

For positive \(S,K,\sigma,\tau\), define
\begin{equation}
 d_1=\frac{\log(S/K)+(r-q+\sigma^2/2)\tau}{\sigma\sqrt\tau},
 \qquad d_2=d_1-\sigma\sqrt\tau.
\end{equation}
Writing the call payoff as \(S_T\mathbf1_{\{S_T>K\}}-K\mathbf1_{\{S_T>K\}}\) separates the expected asset receipt from the strike payment. Evaluating the two lognormal integrals yields
\begin{align}
 C&=Se^{-q\tau}\Phi(d_1)-Ke^{-r\tau}\Phi(d_2),\\
 P&=Ke^{-r\tau}\Phi(-d_2)-Se^{-q\tau}\Phi(-d_1),
 \label{eq:s3_bs}
\end{align}
where \(\Phi\) denotes the standard normal cumulative distribution function \cite{black1973,merton1973}. This is a reference calculation, not a competing numerical PDE solver. It also remains an efficiency baseline for repeated single-asset vanilla pricing.

The implementation will evaluate expiry payoffs directly when \(\tau=0\). At zero volatility, the prices are \(C=(Se^{-q\tau}-Ke^{-r\tau})^+\) and \(P=(Ke^{-r\tau}-Se^{-q\tau})^+\). At \(S=0\), the call is zero and the put is \(Ke^{-r\tau}\). These branches avoid division by zero and undefined logarithms. Reference checks include known prices, limiting cases, and put--call parity:
\begin{equation}
 C-P=Se^{-q\tau}-Ke^{-r\tau}.
\end{equation}

\subsubsection{Reference Greeks}

Delta, \(\Delta=\partial V/\partial S\), measures sensitivity to the underlying price, and Gamma, \(\Gamma=\partial^2V/\partial S^2\), measures how Delta changes with that price. Subscripts \(C\) and \(P\) denote call and put options. With \(\phi\) denoting the standard normal density, the principal sensitivities are
\begin{align}
 \Delta_C&=e^{-q\tau}\Phi(d_1), &
 \Delta_P&=e^{-q\tau}(\Phi(d_1)-1),\\
 \Gamma_C=\Gamma_P&=\frac{e^{-q\tau}\phi(d_1)}{S\sigma\sqrt\tau}, &
 \mathcal V_C=\mathcal V_P&=Se^{-q\tau}\phi(d_1)\sqrt\tau.
\end{align}
Here \(\mathcal V=\partial V/\partial\sigma\) is Vega. Calendar-time Theta is \(\Theta=\partial V/\partial t=-\partial V/\partial\tau\), giving
\begin{align}
 \Theta_C&=-\frac{Se^{-q\tau}\phi(d_1)\sigma}{2\sqrt\tau}
       +qSe^{-q\tau}\Phi(d_1)-rKe^{-r\tau}\Phi(d_2),\\
 \Theta_P&=-\frac{Se^{-q\tau}\phi(d_1)\sigma}{2\sqrt\tau}
       -qSe^{-q\tau}\Phi(-d_1)+rKe^{-r\tau}\Phi(-d_2).
\end{align}
The interest-rate sensitivities, \(\mathrm{Rho}=\partial V/\partial r\), are
\begin{equation}
 \mathrm{Rho}_C=K\tau e^{-r\tau}\Phi(d_2),\qquad
 \mathrm{Rho}_P=-K\tau e^{-r\tau}\Phi(-d_2).
\end{equation}
Vega and Rho are measured per unit change in their decimal inputs; values per percentage point are obtained by multiplying by \(0.01\). Theta is measured per year. Analytical Greeks will be checked against numerical derivatives at smooth interior points. Undefined derivatives at expiry kinks will not be assigned arbitrary finite reference values.

\subsection{Monte Carlo estimation and uncertainty}

Generate \(M\) independent normal samples and terminal prices from Equation~\eqref{eq:s3_terminal}. For discounted payoffs \(Y_m=e^{-r\tau}g(S_T^{(m)})\),
\begin{equation}
 \widehat V_M=\frac1M\sum_{m=1}^M Y_m,\qquad
 s_Y^2=\frac1{M-1}\sum_{m=1}^M(Y_m-\widehat V_M)^2.
\end{equation}
For finite payoff variance, independent sampling gives an unbiased price estimator with variance \(\operatorname{Var}(Y)/M\). The estimated standard error and approximate large-sample 95\% interval are
\begin{equation}
 \widehat{\mathrm{SE}}=\frac{s_Y}{\sqrt M},\qquad
 \widehat V_M\pm1.96\widehat{\mathrm{SE}}.
\end{equation}
The root-mean-square sampling error decreases at order \(M^{-1/2}\) \cite{glasserman2004}. This does not guarantee that one run improves monotonically when samples are added. Convergence will be assessed over independent repetitions, and interval coverage will be measured over runs rather than requiring every interval to contain the analytical value.

The implementation will validate inputs, initialise a recorded seed, generate terminal values, discount the payoffs, and return the price with its uncertainty and sample count. Large simulations will use batches and stable running moment calculations. Tests will check \(\mathbb E^Q[S_T]=Se^{(r-q)\tau}\) and agreement with analytical prices within sampling uncertainty. Runtime will include random-number generation and payoff evaluation. Although the sampling exponent does not explicitly depend on asset dimension, the cost per sample and payoff variance can increase with dimension.

\subsection{Variance reduction and randomised quasi-Monte Carlo}

\subsubsection{Antithetic and control-variate estimators}

Antithetic sampling evaluates each Gaussian draw with both signs and averages the discounted payoffs:
\begin{equation}
 A_m=\tfrac12\bigl(Y(Z_m)+Y(-Z_m)\bigr),\qquad
 \widehat V_{\mathrm{AV}}=\frac1M\sum_{m=1}^M A_m.
\end{equation}
Negative covariance between paired payoffs reduces variance. Standard errors must use the independent pair averages, rather than treating all \(2M\) payoffs as independent. Plain Monte Carlo and antithetic estimates will be compared at equal payoff-evaluation budgets, and the benefit will be measured for each payoff.

A control variate uses a correlated quantity \(X\) with known expectation:
\begin{equation}
 Y_{\mathrm{CV}}=Y-\beta(X-\mathbb E[X]),\qquad
 \beta^*=\frac{\operatorname{Cov}(Y,X)}{\operatorname{Var}(X)}.
\end{equation}
For a single asset, \(X=e^{-r\tau}S_T\) has expectation \(Se^{-q\tau}\). For arithmetic baskets, the discounted geometric basket payoff provides a control. The coefficient will be estimated using an independent pilot sample and fixed for the production sample, keeping pilot costs in the timing budget. Standard errors will use the adjusted observations. Variance-reduction factors will be reported alongside cost, since a lower variance need not imply greater efficiency.

\subsubsection{Scrambled Sobol sampling}

Randomised quasi-Monte Carlo replaces independent uniform draws with scrambled Sobol points \(\mathbf u_m\). The transformation \(z_{m,j}=\Phi^{-1}(u_{m,j})\) produces inputs for terminal simulation. The sampling dimension is one for a vanilla option and \(N\) for an \(N\)-asset terminal basket payoff. Each scramble will use \(M=2^k\) points to preserve the sequence's balance properties. Endpoint handling in the inverse normal transform will be documented.

Uncertainty is calculated from \(R\) independent scrambles, not from the dependent points within each scramble. If their price estimates are \(\widehat V^{(a)}\), then
\begin{equation}
 \overline V=\frac1R\sum_{a=1}^R\widehat V^{(a)},\qquad
 \widehat{\mathrm{SE}}(\overline V)=\frac{s_R}{\sqrt R},
\end{equation}
where \(s_R\) is the sample standard deviation of the replicate estimates. A Student-\(t\) interval will be treated as approximate. Total work is \(RM\) payoff evaluations. The method can improve practical convergence, but its benefit depends on the payoff and effective dimension and will not be assumed in advance \cite{glasserman2004}.

\subsection{Crank--Nicolson finite differences}

\subsubsection{Domain and spatial approximation}

Writing \(U(S,\tau)=V(S,T-\tau)\) gives a forward equation in time to expiry:
\begin{equation}
 U_\tau=\tfrac12\sigma^2S^2U_{SS}+(r-q)SU_S-rU,
 \qquad U(S,0)=g(S).
\end{equation}
The computational domain is \([0,S_{\max}]\). For a call, \(U(0,\tau)=0\) and \(U(S_{\max},\tau)\approx S_{\max}e^{-q\tau}-Ke^{-r\tau}\). Put boundaries are \(Ke^{-r\tau}\) and approximately zero, respectively. The upper boundaries are asymptotic, so domain enlargement will test their influence on the evaluation region.

Divide the asset-price domain into \(J\) equal intervals of width \(h=S_{\max}/J\), giving grid points \(S_j=jh\) for \(j=0,\ldots,J\). Here \(S_{\max}\) is the chosen upper asset-price boundary, and \(U_j\) approximates \(U(S_j,\tau)\) at a fixed time level. Subscripts \(S\), \(SS\), and \(\tau\) denote partial derivatives. Central differences give
\begin{equation}
 U_S\approx\frac{U_{j+1}-U_{j-1}}{2h},\qquad
 U_{SS}\approx\frac{U_{j+1}-2U_j+U_{j-1}}{h^2}.
\end{equation}
Thus the discrete spatial operator has coefficients
\begin{align}
 a_j&=\frac{\sigma^2S_j^2}{2h^2}-\frac{(r-q)S_j}{2h},\\
 b_j&=-\frac{\sigma^2S_j^2}{h^2}-r,\\
 c_j&=\frac{\sigma^2S_j^2}{2h^2}+\frac{(r-q)S_j}{2h}.
\end{align}
For the interior unknowns, this gives \(\mathbf u'=A\mathbf u+\mathbf f(\tau)\), with tridiagonal \(A\). The first and last components of \(\mathbf f\) contain the lower and upper boundary contributions.

\subsubsection{Time stepping and implementation}

Divide the time-to-expiry interval \([0,\tau_{\max}]\) into \(N_t\) steps of size \(\Delta\tau=\tau_{\max}/N_t\), where \(\tau_{\max}\) is the longest maturity being solved. At time level \(\tau_n=n\Delta\tau\), let \(\mathbf u^n\) contain the interior grid values. With \(I\) denoting the identity matrix, the Crank--Nicolson update is
\begin{equation}
 (I-\tfrac{\Delta\tau}{2}A)\mathbf u^{n+1}
 =(I+\tfrac{\Delta\tau}{2}A)\mathbf u^n
 +\tfrac{\Delta\tau}{2}\bigl(\mathbf f(\tau_n)+\mathbf f(\tau_{n+1})\bigr).
 \label{eq:s3_cn}
\end{equation}
The method is second order in space and time for suitably smooth, resolved solutions. The payoff kink can cause oscillations and degrade Greek convergence. The first two full time intervals will therefore be replaced by four backward-Euler half-steps, followed by Crank--Nicolson. This Rannacher damping suppresses initial oscillations without treating an undamped implementation as the strongest classical baseline \cite{pooley2003}.

Implementation consists of grid construction, payoff initialisation, operator assembly, damped starting steps, and repeated solution of Equation~\eqref{eq:s3_cn}. Boundary values will be updated at each time level. A banded or tridiagonal solver will be used; a dense matrix inverse will not be formed. With \(J\) spatial intervals and \(N_t\) steps, work is \(O(JN_t)\). Working memory is \(O(J)\), or \(O(JN_t)\) if the complete surface is retained for repeated queries.

Spatial and temporal refinement will be tested separately while keeping the other error sources small. An observed spatial order is \(p=\log(E(h)/E(h/2))/\log2\), where \(E(h)\) is the error measured using the same norm and evaluation points for grid spacing \(h\). Domain truncation and off-grid interpolation will receive separate checks. Stability alone does not establish accuracy or monotonicity. Low-volatility cases can also make centred convection unsuitable; any change to the spatial stencil will be identified explicitly and its order rechecked.

\subsection{Numerical Greeks and correlated baskets}

\subsubsection{Sensitivity estimators}

For a price approximation \(\widehat V\) and positive spot-price perturbation \(\varepsilon\), central differences yield
\begin{align}
 \widehat\Delta&=\frac{\widehat V(S+\varepsilon)-\widehat V(S-\varepsilon)}{2\varepsilon},\\
 \widehat\Gamma&=\frac{\widehat V(S+\varepsilon)-2\widehat V(S)+\widehat V(S-\varepsilon)}{\varepsilon^2}.
\end{align}
Crank--Nicolson will use grid neighbours for these derivatives. Off-grid interpolation will be documented; differentiating a piecewise-linear interpolant twice is not a suitable Gamma estimator. Monte Carlo bumps will share random samples, with standard errors estimated from the per-sample combined differences. Several perturbation sizes will be used because reducing a bump can amplify solver error, sampling noise, and floating-point cancellation. Vega and Rho use parameter perturbations, and Theta uses the negative maturity derivative, with one-sided treatment where necessary.

A call also admits the pathwise estimator
\begin{equation}
 \widehat\Delta_{\mathrm{PW}}=\frac{e^{-r\tau}}{M}
 \sum_{m=1}^M\mathbf1_{\{S_T^{(m)}>K\}}\frac{S_T^{(m)}}S.
\end{equation}
It provides an additional Delta comparison \cite{glasserman2004}. Naively differentiating the indicator again misses the payoff boundary contribution and does not give a valid ordinary pathwise Gamma estimator. The initial Gamma study will therefore use common-random-number finite differences with both bias and uncertainty checks.

\subsubsection{Basket simulation and analytical control}

For baskets of \(N=2,3,5\) assets, let \(\rho_{ij}\) be the correlation between the Brownian increments driving assets \(i\) and \(j\). The positive-definite correlation matrix \(R=(\rho_{ij})\) is factored as \(R=LL^{\mathsf T}\). Here \(L\) is a lower-triangular Cholesky factor. A vector \(\mathbf Z\) of \(N\) independent standard normal variables gives correlated inputs \(\mathbf X=L\mathbf Z\). Each asset uses Equation~\eqref{eq:s3_terminal} with its own current price \(S_i\), dividend yield \(q_i\), volatility \(\sigma_i\), and correlated normal draw \(X_i\). With samples stored as rows, the array operation is \(ZL^{\mathsf T}\). Symmetry, unit diagonal, and positive definiteness will be checked before Cholesky factorisation; singular valid cases require a suitable alternative factorisation.

Let \(S_{i,T}\) denote asset \(i\)\textquotesingle s terminal price and \(w_i\) its fixed basket weight. The arithmetic basket is \(B_T=\sum_iw_iS_{i,T}\), with \(w_i\geq0\) and \(\sum_iw_i=1\). Pricing averages its discounted payoff \((B_T-K)^+\). For the geometric basket \(G_T=\prod_iS_{i,T}^{w_i}\), \(\log G_T\) is normally distributed with mean \(m_G\) and variance \(s_G^2\), given by
\begin{align}
 m_G&=\sum_iw_i\log S_i+
 \left(r-\sum_iw_iq_i-\tfrac12\sum_iw_i\sigma_i^2\right)\tau,\\
 s_G^2&=\tau\sum_i\sum_jw_iw_j\sigma_i\sigma_j\rho_{ij}.
\end{align}
For \(s_G>0\), let \(d_{2,G}=(m_G-\log K)/s_G\) and \(d_{1,G}=d_{2,G}+s_G\). The call reference is
\begin{equation}
 C_G=e^{-r\tau}\left[e^{m_G+s_G^2/2}\Phi(d_{1,G})-K\Phi(d_{2,G})\right].
\end{equation}
This supplies validation and the known expectation for a geometric-payoff control variate. The arithmetic basket is not generally lognormal, so its independent reference estimates retain statistical uncertainty. The \(N=1\) reduction and the generated Gaussian covariance will be checked.

Basket sensitivities will include asset-wise Delta, diagonal Gamma, and selected cross-Gammas. Here \(\widehat V_{ab}\) denotes the estimated price with \(S_i\) shifted by \(a\varepsilon_i\) and \(S_j\) by \(b\varepsilon_j\), where \(a,b\in\{-1,+1\}\) and \(\varepsilon_i,\varepsilon_j>0\) are the spot perturbation sizes. Perturbing two asset prices gives
\begin{equation}
 \widehat\Gamma_{ij}=
 \frac{\widehat V_{++}-\widehat V_{+-}-\widehat V_{-+}+\widehat V_{--}}
 {4\varepsilon_i\varepsilon_j},\qquad i\ne j,
\end{equation}
using the same Gaussian samples in the four valuations. Correlation and all other parameters remain fixed during these spot perturbations.

\subsection{Validation and comparison protocol}

Each routine will accept explicit financial inputs and return its price, available sensitivities, and diagnostics. Separate experiment scripts will record seeds, grid or sample sizes, software versions, and output tables. Validation will proceed from analytical identities to stochastic estimates and grid convergence, then to basket cases.

Prices and each Greek will be compared through root-mean-square error (RMSE) and maximum absolute error on a common independent evaluation set. Relative errors will be restricted to reference values away from zero. Financial checks will include applicable price bounds, parity, monotonicity, and convexity, with statistical uncertainty considered for simulation estimates.

Cost will include all work needed for the requested output, including variance-reduction pilots. A single price, a complete surface, and Greeks are distinct workloads. Repeated-query comparisons will account for interpolation from an existing classical surface, rather than forcing a full recalculation for every query. Runtime, memory, numerical error, and uncertainty will be reported together. For fixed-model benchmarks, differences from the analytical or high-accuracy numerical reference measure solver error; differences from actual option quotes additionally involve calibration, model assumptions, and data quality.




% --------------------------------------------------
% References
% --------------------------------------------------
\begin{thebibliography}{99}

\bibitem{black1973}
F. Black and M. Scholes,
``The pricing of options and corporate liabilities,''
\textit{Journal of Political Economy},
vol.~81, no.~3, pp.~637--654, 1973.
\url{https://doi.org/10.1086/260062}

\bibitem{merton1973}
R. C. Merton,
``Theory of rational option pricing,''
\textit{The Bell Journal of Economics and Management Science},
vol.~4, no.~1, pp.~141--183, 1973.
\url{https://doi.org/10.2307/3003143}

\bibitem{glasserman2004}
P. Glasserman,
\textit{Monte Carlo Methods in Financial Engineering}.
New York: Springer, 2004.
\url{https://doi.org/10.1007/978-0-387-21617-1}

\bibitem{raissi2019}
M. Raissi, P. Perdikaris, and G. E. Karniadakis,
``Physics-informed neural networks: A deep learning framework
for solving forward and inverse problems involving nonlinear
partial differential equations,''
\textit{Journal of Computational Physics},
vol.~378, pp.~686--707, 2019.
\url{https://doi.org/10.1016/j.jcp.2018.10.045}

\bibitem{lu2021}
L. Lu, P. Jin, G. Pang, Z. Zhang, and G. E. Karniadakis,
``Learning nonlinear operators via DeepONet based on the
universal approximation theorem of operators,''
\textit{Nature Machine Intelligence},
vol.~3, pp.~218--229, 2021.
\url{https://doi.org/10.1038/s42256-021-00302-5}

\bibitem{wang2021deeponet}
S. Wang, H. Wang, and P. Perdikaris,
``Learning the solution operator of parametric partial
differential equations with physics-informed DeepONets,''
\textit{Science Advances},
vol.~7, no.~40, article eabi8605, 2021.
\url{https://doi.org/10.1126/sciadv.abi8605}

\bibitem{krishnapriyan2021}
A. S. Krishnapriyan, A. Gholami, S. Zhe,
R. M. Kirby, and M. W. Mahoney,
``Characterizing possible failure modes in physics-informed
neural networks,''
\textit{Advances in Neural Information Processing Systems},
vol.~34, pp.~26548--26560, 2021.
\url{https://arxiv.org/abs/2109.01050}

\bibitem{wang2021gradients}
S. Wang, Y. Teng, and P. Perdikaris,
``Understanding and mitigating gradient flow pathologies
in physics-informed neural networks,''
\textit{SIAM Journal on Scientific Computing},
vol.~43, no.~5, pp.~A3055--A3081, 2021.
\url{https://doi.org/10.1137/20M1318043}

\bibitem{pooley2003}
D. M. Pooley, K. R. Vetzal, and P. A. Forsyth,
``Convergence remedies for non-smooth payoffs in option pricing,''
\textit{Journal of Computational Finance},
vol.~6, no.~4, pp.~25--40, 2003.
\url{https://cs.uwaterloo.ca/~paforsyt/report.pdf}

\bibitem{nseData}
National Stock Exchange of India,
``Historical reports and equity derivatives reports.''
Online data resources.
\url{https://www.nseindia.com/static/resources/historical-reports-capital-market-daily-monthly-archives}
and
\url{https://www.nseindia.com/all-reports-derivatives}

\end{thebibliography}

\end{document}